\subsection{Explicit incidence laws in low dimensions}
The distinction between oriented elements and orientation orbits is
already necessary for consistency. If an irreflexive relation on a
nonempty set with free reversal is invariant under reversing either
argument, it cannot be ordinary inequality. Indeed
$x\ne\bar x$ would give $x\perp\bar x$, hence
$\bar x\perp\bar x$. The relation
$x\perp y\Longleftrightarrow q(x)\ne q(y)$ is a consistent
alternative. It distinguishes underlying cells, not their two
orientations.

We state the low-dimensional incidence conditions explicitly,
writing $\triangleleft$ for immediate incidence.
Let the orbit poset be graded in nonnegative dimensions, let immediate
incidences project exactly to covers, and require simultaneous
reversal. Choose positive vertex orientations. Impose these two laws:
\begin{enumerate}
\item Each oriented edge has exactly one positive and exactly one
negative vertex immediately below it.
\item If $z<x$ has dimension difference two, then for each orientation
$Z$ of $z$ and each orientation $X$ of $x$ there is exactly one
oriented intermediate cell $Y$ with
$Z\mathrel{\triangleleft}Y\mathrel{\triangleleft}X$.
\end{enumerate}
Unlike Theorem~\ref{thm:cw}, these two laws do not prohibit both
orientations of a facet from being incident to the same oriented cell.
That distinction is essential in dimension two.

\begin{proposition}[The graph and boundary dictionary]\label{prop:lowdim-dictionary}
In dimensions at most one, the displayed incidence data are exactly
graphs specified by a vertex set $V$, an oriented-edge set $E$ with
free reversal, and an origin map $s:E\to V$. The terminus is
$t(e)=s(\bar e)$. Loops and parallel edges are allowed.

In dimension at most two, the oriented edges below each chosen
two-cell orientation form a disjoint union of finite directed cycles
and two-sided infinite directed paths. Every incident vertex occurs
once as origin and once as terminus. Conversely any such boundary
data determine incidence data satisfying the two laws. If each cell
has finitely many faces, only finite cycles occur.
\end{proposition}
\begin{proof}
The first law gives the unique origin $s(e)$. Reversal gives the
unique negative terminal vertex $t(e)=s(\bar e)$. Thus the signed
incidences are exactly
\[
 s(e)^+\triangleleft e,\qquad t(e)^-\triangleleft e,
\]
and their simultaneous reversals. These formulas also reconstruct
the incidence relation from any graph of the stated kind.

Fix a two-cell orientation $f$ and let $B_f$ be its incident
oriented-edge set. Applying the second law to the positive and
negative orientations of a boundary vertex gives exactly one
outgoing and one incoming edge in $B_f$. Following the outgoing
edge is therefore a bijection of the boundary vertex set.
Its orbits are finite cycles or free $\Z$-orbits, the latter
being two-sided infinite paths. Conversely these conditions give
exactly one two-step path for each endpoint orientation, so they
recover the second law. A finite face set excludes an infinite orbit.
\end{proof}
For example, the edges $e_i:v_i\to v_{i+1}$, $i\in\Z$, may
all be placed below one formal positive two-cell. Every incidence
chain has length at most two, yet its boundary is a line.
Finite dimension does not imply finite cell closure.

There is also a substantive difference between matching and
unmatched signs in an incidence-lifting condition. Write $D_X$ for
the immediate predecessor set of $X$. A condition that, under
identity refinement, transfers an incidence into either $X$ or
$\bar X$ to an incidence into $X$ forces
$D_X=D_{\bar X}$. With the displayed laws this leaves exactly
graphs consisting of loops and isolated vertices. To see this,
simultaneous reversal makes every $D_X$ reversal-stable. Any
two-step chain then has both $Y$ and $\bar Y$ as intermediates,
contrary to uniqueness. Thus no cell of dimension at least two
remains. At an edge, reversal stability and the unique endpoint
law make its positive origin and negative terminus lie over the
same vertex. Conversely loop graphs satisfy all these conditions.
If the conclusion uses the same upper sign as the hypothesis,
the lifting condition under identity refinement is tautological.

\begin{proposition}[The exact regular class in dimension two]
\label{prop:lowdim-regular}
For the preceding incidence data in dimensions at most two, assume
each cell has finitely many faces. They arise from a regular CW
complex exactly when no edge is a loop, every two-cell boundary
is nonempty and connected, and no such boundary uses both
orientations of one underlying edge. No global bound on the
number of cells or on vertex valence is required.
\end{proposition}
\begin{proof}
A finite connected directed boundary is one cycle, by
Proposition~\ref{prop:lowdim-dictionary}. Its vertex set consists
of the distinct elements in one orbit. The two exclusions make
its underlying edges distinct and rule out a one-edge loop.
A two-edge cycle with distinct parallel edges is allowed.
Thus each prescribed boundary embeds as a circle subcomplex of
the underlying graph. Realize non-loop edges by closed intervals
and attach a disk to each boundary circle by a homeomorphism.
The resulting closure-finite weak cell structure is a regular
CW complex. Conversely a regular edge is an embedded closed
interval and a regular two-cell has an embedded boundary circle.
It has neither a loop edge nor a repeated underlying edge, and
its boundary is connected. These give precisely the conditions.
\end{proof}

Excluding loops alone is insufficient. Take two vertices $a,b$,
one non-loop edge $e$ from $a$ to $b$, and one formal two-cell
whose positive boundary contains both $e$ and $\bar e$.
Reverse all incidences simultaneously for the negative face.
Each boundary vertex has one incoming and one outgoing edge,
so both incidence laws hold. The unoriented boundary is a closed
interval. It cannot be the boundary circle of a regular two-cell.
This example has four unoriented cells and eight oriented cells.

\subsection{A finite obstruction with all vertex links connected}
\begin{counterexample}\label{ex:torus-55}
There is a signed incidence system with $55$ unoriented cells,
of maximal dimension three, such that every signed diamond holds,
every boundary of positive dimension is connected, and every
vertex link in its order complex is connected, but it is not
the face poset of a regular CW complex.
\end{counterexample}
\begin{proof}
Triangulate the plane by unit squares split along their increasing
diagonals, and quotient by $3\Z^2$. The vertex set is
$V=(\Z/3)^2$. With $u=(1,0)$ and $v=(0,1)$, the two positively
oriented triangles based at $w\in V$ are
\[
 (w,w+u,w+u+v),\qquad (w,w+u+v,w+v).
\]
Every triangle embeds in the quotient because its coordinate span
is smaller than three. The two triangle types have distinct
translation orbits, each of size nine. The three edge directions
$u,v,u+v$ similarly give three orbits of nine edges. Thus this
torus triangulation has $9$ vertices, $27$ edges and $18$ triangles.
Adjoin a formal greatest element $\Omega$ of dimension three to
its nonempty face poset, obtaining $55$ cells.

Orient the edges arbitrarily and the triangles by the displayed
orientations. Give the incidence of each triangle into $\Omega$
coefficient $+1$. Every edge lies in two triangles with opposite
induced signs. Below a triangle, the ordinary simplex boundary
has the same cancellation at every vertex. These give every
rank-two diamond, together with the opposite endpoint signs on
edges. Doubling each cell by its two orientations gives the
required signed incidence system.

The two-cell boundaries are circles and the boundary of $\Omega$
is the connected torus. The link of the vertex $\Omega$ in the
order complex is the subdivided torus, hence connected. At an
old cell $c$, its vertex link is
\[
 \Delta(P_{<c})*\Delta(P_{>c}),
\]
where $P$ now includes $\Omega$. The upper interval has maximum
$\Omega$, so its order complex is a nonempty cone. This makes
the link connected; for a minimal cell the lower complex is
empty and the link is just the upper cone.

Finally the strict lower complex of $\Omega$ is a torus and has
first rational Betti number two, while a two-sphere has first
Betti number zero. Thus it cannot be the boundary of a regular
three-cell. The accompanying exact finite calculation gives
boundary ranks $8$ and $17$, hence Betti numbers $(1,2,1)$,
and verifies all $55$ vertex-link connectedness tests.
\end{proof}
The global links in this example must be distinguished from
the lower intervals appearing in Theorem~\ref{thm:cw}.
